% !PRINTED 207
\weberpage{207}

1.~\gesperrt{Jedem Systeme conjugirter Pole der Gruppe $G$
entspricht eine invariante Form, deren Grad
gleich der Anzahl der Pole des Systemes ist, und
deren Wurzeln eben diese Pole sind.}

Wir bezeichnen mit $f_1, f_2, \ldots$ die zu den verschiedenen
Systemen conjugirter Pole gehörigen invarianten Formen, die
von den Graden $\mu, \mu', \ldots$ sind und die wir die \gesperrt{Grundformen}
der Gruppe nennen wollen.

Ist dann $F(x_1, x_2)$ irgend eine invariante Form der
Gruppe $G$ und $f(x, y)$ eine Grundform, und hat $F(x, 1) = 0$
mit $f(x, 1) = 0$ eine gemeinsame Wurzel, so müssen alle Wurzeln
 von $f = 0$ zugleich Wurzeln von $F = 0$ sein. Denn
nach Voraussetzung haben die beiden Functionen $F(x, 1)$ und
$F[\psi(x), 1]$ dieselben Wurzeln. Wenn also $F(a, 1) = 0$ ist,
so ist auch $F[\psi(a), 1] = F(a_1, 1) = 0$. Wir haben also den
folgenden Satz:

2.~\gesperrt{Hat eine zu $G$ gehörige invariante Form $F(x_1, x_2)$
mit einer der Grundformen $f(x_1, x_2)$ einen Theiler
gemein, so ist $F(x_1, x_2)$ durch $f(x_1, x_2)$ theilbar.}

Ist $F(x_1, x_2)$ eine invariante Form der Gruppe $G$, und $\xi$
eine Wurzel der Gleichung $F(x, 1) = 0$, so sind auch, wenn
$\psi$ die Elemente der Gruppe $G$ durchläuft, die sämmtlichen
$n$ Grössen $\psi(\xi)$ Wurzeln derselben Gleichung. Wenn also der
Grad von $F$ niedriger ist, als der Grad der Gruppe, so können
diese Grössen nicht alle von einander verschieden sein, und es
folgt für irgend zwei von einander verschiedene Substitutionen
$\chi, \psi$ von $G$
\[
\chi(\xi) = \psi(\xi)
\]
oder
\[
\xi = \chi^{-1}\psi(\xi),
\]
d. h. $\xi$ muss unter den Polen der Gruppe $G$ vorkommen, und
es ist also $F$ nach dem Satze 2. durch eine der Grundformen
theilbar. Da wir auf den Quotienten der Division dieselbe Schlussweise
 anwenden können, so folgt:

3.~\gesperrt{Eine invariante Form $F$ der Gruppe $G$, deren
Grad niedriger ist, als der Grad der Gruppe,
ist, von einem constanten Factor abgesehen, ein
Product von Potenzen der Grundformen.}
%CONT>>
